% Einheit 4 von 5 – Dezimalzahlen (bank/brueche-dezimalzahlen/e4.jsonl, zusammenbau v0.1)
%% TODO Zweigzeile Teil 1: was hier gelernt wird (Satz in Schülersprache aus dem Einheitstitel) – Einheit 4
\einheitenkopf[e4]{Einheit 4 von 5 · Dezimalzahlen}
\zweigzeile{ab Kl. 5 · P10 oft}
\setcounter{aufgabe}{20}

%% TODO Titel: Ich-kann-Satz fehlt in der Bank (Platzhalter: Kettenname) – Nr. 21
%% TODO Anweisung: ein Satz über der ersten Teilaufgabe fehlt in der Bank – Nr. 21
\begin{aufgabe}{Umwandeln}
%% TODO \rechenplatz{n}: Schrittzahl des Grundfalls fehlt in der Bank (nur bei fester Schreibform, 2.3 b)
\begin{teile}
\teil $\frac{9}{100}$ – wie viele Stellen nach dem Komma? \leerfeld
\teil $\frac{6}{10}$ – als Dezimalzahl? \leerfeld
\teil $\frac{7}{100}$ – als Dezimalzahl? \leerfeld
\teil Trage $0{,}4$ und $1{,}3$ am Zahlenstrahl ein.

\zahlenstrahl[xmin=0,xmax=2,xstep=0.1,karo=0.6]{0/0, 1/1, 2/2}
\teil $\frac{2}{5}$ – als Dezimalzahl? \leerfeld
\teil $0{,}6$ – als gekürzter Bruch? \leerfeld
\teil $\frac{7}{8}$ – als Dezimalzahl? Teile den Zähler durch den Nenner. \leerfeld
\teil $\frac{4}{9}$ – als Dezimalzahl? Die Division geht nicht auf; schreibe mit Periodenstrich. \leerfeld
\teil Welche Aussage ist wahr? Kreuze an. \hfill (P10 2015 OS)\\ \kreuz{$2{,}5 < \frac{5}{2}$}\\ \kreuz{$\frac{13}{5} > \frac{5}{2}$}\\ \kreuz{$\sqrt{5} > \frac{5}{2}$}
\end{teile}
\end{aufgabe}

%% TODO Titel: Ich-kann-Satz fehlt in der Bank (Platzhalter: Kettenname) – Nr. 22
%% TODO Anweisung: ein Satz über der ersten Teilaufgabe fehlt in der Bank – Nr. 22
\begin{aufgabe}{Zahl in die Stellenwerttafel eintragen und ablesen (Zehntel, Hundertstel, Tausendstel)}
\begin{teile}
\teil Welche Dezimalzahl steht in der Stellenwerttafel?

\sachtabelle{cccc}{Einer & Zehntel & Hundertstel & Tausendstel}{2 & 0 & 4 & 7}
\leerfeld
\end{teile}
\end{aufgabe}

%% TODO Titel: Ich-kann-Satz fehlt in der Bank (Platzhalter: Kettenname) – Nr. 23
%% TODO Anweisung: ein Satz über der ersten Teilaufgabe fehlt in der Bank – Nr. 23
\begin{aufgabe}{Nachbarzahlen (Nachbar-Einer, -Zehntel, -Hundertstel) und Zählen in Schritten (0,2er-Schritte; über 2,9 hinaus)}
\begin{teile}
\teil Nenne die Nachbarzehntel von $4{,}37$. \leerfeld und \leerfeld
\end{teile}
\end{aufgabe}

%% TODO Titel: Ich-kann-Satz fehlt in der Bank (Platzhalter: Kettenname) – Nr. 24
%% TODO Anweisung: ein Satz über der ersten Teilaufgabe fehlt in der Bank – Nr. 24
\begin{aufgabe}{Umwandeln – Fehler finden, Begründen, Darstellungswechsel, Anwendung}
\begin{teile}
\teil Jana schreibt: $\frac{1}{4} = 1{,}4$. Finde den Fehler.
\teil Warum ist $\frac{9}{20} = 0{,}45$? Begründe.
\teil Welche Dezimalzahl zeigt der gefärbte Teil des Streifens?

\bruchrechteck{2}{10}
\leerfeld
\teil Ein Rezept verlangt $\frac{3}{5}$ l Milch. Der Messbecher hat eine Skala in Litern mit Komma. Bis zu welcher Marke füllst du? \leerfeld[l]
\end{teile}
\end{aufgabe}
